% le caratteristiche richieste dall'università sono ellamport1986latexencate qui: https://stem.elearning.unipd.it/mod/book/view.php?id=234&chapterid=46#modalita
% 12pt: font richiesto dall'università
% twoside: i margini interni ed esterni sono scambiati per le pagine "a sinistra" e "a destra"
% openright: i capitoli cominciano in pagine dispari ("a destra")
% extreport: supporta 12pt
\documentclass[12pt,a4paper,twoside,openright]{extreport}

\usepackage{amsmath}                            % per avere più controllo sulle equazioni 
\usepackage{csquotes}                           % per le citazioni
\usepackage{enumitem}                           % per avere più controllo sulle enumerazioni
\usepackage[
    a4paper,
    top=2cm,bottom=2cm,
    outer=2cm,inner=3cm,
    includeheadfoot
]{geometry}                                     % margini richiesti dall'università
\usepackage{graphicx}                           % per le immagini
\usepackage{icomma}                             % per separare le cifre decimali con una virgola
\usepackage{minted}                             % per il codice con la colorazione della sintassi
\usepackage[a-1a]{pdfx}                         % formato richiesto dall'università
\usepackage[output-decimal-marker={,}]{siunitx} % per le unità di misura
\usepackage{subcaption}                         % per le sottodidascalie
\usepackage{svg}                                % per inserire file SVG

\usepackage[italian]{babel}
\selectlanguage{italian}

\usepackage[backend=bibtex,style=ieee]{biblatex}
\bibliography{bibliografia}

\usepackage{fontspec}
\setmainfont{Times New Roman} % carattere richiesto dall'università

\usepackage{setspace}
\onehalfspacing % interlinea richiesta dall'università

\sloppy % per evitare che il testo in \verb finisca oltre i margini

% questi valori vengono usati nella composizione del frontespizio
\title{Titolo della tesi}
\author{Stella Filippo}
\date{GG/MM/AAAA}
\newcommand{\supervisor}{Prof. Cagnazzo Marco}
\newcommand{\assistantsupervisor}{Cognome Nome}

\begin{document}
    \pagenumbering{roman}
    \pagestyle{empty} % per le prime pagine, non mostrare il numero di pagina

    \begin{titlepage}
    % solo il frontespizio deve essere simmetrico rispetto ai margini interno ed esterno
    \newgeometry{hmargin=2.5cm,vmargin=2cm}
        \begin{figure}
            \centering
            \begin{subfigure}[b]{0.4\textwidth}
                \includegraphics[width=\textwidth]{Immagini/logo_unipd}
            \end{subfigure}
            \hfill
            \begin{subfigure}[b]{0.3\textwidth}
                \includegraphics[width=\textwidth]{Immagini/logo_dei}
            \end{subfigure}
        \end{figure}
    
        \vspace*{\stretch{0.5}}
    
        \begin{center}
            \makeatletter % serve per poter usare \@...

            % NOTA: il Times New Roman non supporta il maiuscoletto.
            \textsc{DIPARTIMENTO DI INGEGNERIA DELL'INFORMAZIONE}\\
            \vspace*{\stretch{0.1}}
            \textsc{CORSO DI LAUREA IN INGEGNERIA INFORMATICA}
    
            \vspace*{\stretch{0.5}}
            \LARGE
            \textbf{\@title}
    
            \vspace*{\stretch{1}}
            \normalsize
            \begin{tabular*}{\textwidth}{l @{\extracolsep{\fill}} r}
                \textbf{Relatore} & \textbf{Laureando} \\
                \supervisor       & \@author           \\
                \\
%                \textbf{Correlatore} \\
%                \assistantsupervisor \\
            \end{tabular*}
    
            \vspace*{\stretch{2}}
            \textsc{ANNO ACCADEMICO 2022-2023} \\
            \vspace*{\stretch{0.1}}
            Data di laurea \@date
        
            \makeatother % serve dopo \makeatletter
        \end{center}
    \restoregeometry
\end{titlepage}

    \cleardoublepage
    
    \vspace*{\stretch{1}}
\begin{flushright}
    \textit{Apes. Together. Strong.}
\end{flushright}
Testo di prova
\vspace{\stretch{4}}
    \cleardoublepage

    \pagestyle{plain} % comincia a mostrare il numero di pagina

    % NOTA: l'ambiente \abstract rimuove il numero della pagina e resetta il contatore delle pagine. 
    \chapter*{Sommario}
    Descrizione dell'obiettivo della tesi.
    \cleardoublepage

    \tableofcontents
    \cleardoublepage
    
    \listoffigures
    \cleardoublepage % per assicurarsi che la numerazione araba cominci col primo capitolo
    
    \pagenumbering{arabic}

    \chapter{Capitolo di prova}

Questo capitolo può servire come riferimento per alcuni comandi utili di \LaTeX.

\section{Esempio di sezione}

\subsection{Esempio di sottosezione}

\begin{displayquote}[J. R. R. Tolkien]
    The Lord of the Rings \textit{is one of these things: \\ if you like it, you do; \\ if you don't, then you boo!}
\end{displayquote}

Questa qui sopra è una citazione. Bella, vero?

Questo è un riferimento: \cite{lamport1986latex}.

Questa è un'espressione matematica in riga: $y = f(x)$. Questa è un'espressione matematica in un blocco a sé:

\begin{equation}
    \int\limits_{-\infty}^\infty e^{-x^2}dx = \sqrt{\pi}
\end{equation}

Questo è codice in riga: \verb|for(int i = 0; i < 1; i++);|. Questo è un listato semplice, estratto da \textit{The C programming language} (1988), p. 6:

\begin{verbatim}
    #include <stdio.h>

    main( )
    {
        printf("hello, world\n");
    }
\end{verbatim}

Questo è un listato sbrilluccicoso:

\begin{minted}[mathescape,gobble=2]{python}
    # can your editor do this? $\displaystyle{f(x) = m + \sum\limits_{k=1}^\infty A_k \cos(k x - \varphi_k)}$

    import this

    print('hello, world')
\end{minted}
    \cleardoublepage

    \printbibliography[heading=bibintoc]
\end{document}
